% !TeX root = ../main.tex

\appendix
\ctexset{chapter={name={附录},number={\Alph{chapter}},format+={\zihao{-2}\usefonttimes},afterskip={0.5\baselineskip},beforeskip={0.8\baselineskip}}}
\setcounter{chapter}{0}
\setcounter{table}{0}
\renewcommand{\thetable}{A\arabic{table}}
\chapter{年度收益完整表}
\begin{table}[htbp]
\centering
\caption{年度复利收益（\%）}
\renewcommand{\arraystretch}{1.18}
\footnotesize
\setlength{\tabcolsep}{3pt}
\begin{tabularx}{\textwidth}{@{}p{1.05cm}*{7}{C}@{}}
\toprule
年份 & 月数 & E & P1 & P2 & P3 & P4 & 上证 \\
\midrule
2003 & 9 & 36.32 & 36.32 & 36.32 & 26.16 & 26.16 & -0.90 \\
2004 & 12 & 6.84 & 6.84 & 6.84 & 4.79 & 4.79 & -15.40 \\
2005 & 12 & -0.05 & -0.05 & -0.05 & 0.33 & 0.33 & -8.33 \\
2006 & 12 & 126.39 & 126.39 & 126.39 & 84.78 & 84.78 & 130.43 \\
2007 & 12 & 446.28 & 446.28 & 446.28 & 182.09 & 182.09 & 96.66 \\
2008 & 12 & -65.02 & -60.67 & -52.44 & -25.19 & -20.23 & -65.39 \\
2009 & 12 & 91.75 & 169.65 & 143.31 & 61.08 & 60.70 & 79.98 \\
2010 & 12 & 39.48 & 31.00 & 39.52 & 18.98 & 27.37 & -12.15 \\
2011 & 12 & -9.15 & -17.17 & -20.47 & -18.11 & -20.82 & -21.68 \\
2012 & 12 & 20.58 & 21.73 & -2.02 & 19.08 & -3.07 & 3.17 \\
2013 & 12 & 37.33 & 19.28 & 37.24 & 13.93 & 23.70 & -6.75 \\
2014 & 12 & 49.37 & 50.14 & 53.96 & 51.88 & 46.43 & 52.87 \\
2015 & 12 & 38.87 & 62.56 & 71.35 & 54.29 & 71.35 & 8.49 \\
2016 & 12 & 30.82 & 5.14 & -3.22 & 1.36 & -1.46 & -12.31 \\
2017 & 12 & 5.18 & 14.10 & 6.97 & 13.01 & 6.69 & 6.56 \\
2018 & 12 & -6.41 & -25.06 & -27.19 & -25.06 & -27.19 & -24.59 \\
2019 & 12 & 55.91 & 34.89 & 34.50 & 33.01 & 32.15 & 22.30 \\
2020 & 12 & 17.81 & 25.10 & 24.40 & 25.19 & 24.88 & 13.87 \\
2021 & 12 & -17.86 & 48.86 & 58.92 & 48.82 & 55.89 & 4.80 \\
2022 & 12 & 13.95 & 8.08 & 21.89 & 3.94 & 13.22 & -15.13 \\
2023 & 12 & -8.77 & 9.89 & 4.54 & 9.09 & 1.12 & -3.70 \\
2024 & 12 & 29.52 & 38.52 & 17.09 & 38.12 & 16.90 & 12.67 \\
2025 & 12 & 48.35 & 48.82 & 30.44 & 44.26 & 30.33 & 18.41 \\
2026 & 9 & 15.50 & 9.99 & -7.71 & 8.35 & -7.20 & -3.19 \\
\bottomrule
\end{tabularx}
\par\smallskip
\begin{minipage}{\textwidth}\small 注：B与E在共同区间完全一致，不重复列B的自然年度短表。2003年为4至12月，2026年为1至9月；其余为12个月。均为复利实际收益，不将不足一年区间在此表内年化。\end{minipage}
\end{table}

\setcounter{table}{0}
\renewcommand{\thetable}{B\arabic{table}}
\chapter{历史冻结的信息时点诊断}
本附录仅使用此前已核验的261个月结果，不混入新追加21个月。旧基准月收益与本次E的前261个月逐月一致。固定滞后一季不是完整公告约束；完整约束需要同比分母、TTM依赖和历史成长记录等财务依赖均已公告，并使用当时可得的最新财务期。该计算仍使用最终会计值，不能称为完整还原历史版本的数据库。

\begin{table}[htbp]
\centering
\caption{历史冻结公告敏感性（仅261个月）}
\renewcommand{\arraystretch}{1.18}
\small
\setlength{\tabcolsep}{4.5pt}
\begin{tabularx}{\textwidth}{@{}p{3.15cm}*{3}{C}@{}}
\toprule
实施口径 & 年化收益/\% & 波动/\% & 最大回撤/\% \\
\midrule
报告期映射 & 25.40 & 32.50 & -70.11 \\
固定滞后一季 & 10.38 & 32.22 & -72.30 \\
完整公告约束 & 0.44 & 33.37 & -89.91 \\
仅当期公告 & 6.31 & 32.66 & -72.30 \\
公告覆盖匹配 & 24.85 & 31.78 & -68.06 \\
\bottomrule
\end{tabularx}
\end{table}

\begin{figure}[htbp]
\centering
\includegraphics[width=\linewidth]{fig16}
\caption{历史冻结区间的公告敏感性}
\par\smallskip
\begin{minipage}{0.97\linewidth}\small 注：仅2003-04至2024-12；不是四条优化路径的公告约束收益，不能外推为2026年统计。资料来源：CSMAR，作者计算。\end{minipage}
\end{figure}

覆盖匹配仍允许提前使用财务值，仅用于观察覆盖剔除能否解释差异。完整约束同时改变财务期和排名集合，不能用两者CAGR之差识别单一偏差的因果效应。此前HAC、因子回归和区块自助研究连同原报告保留在归档目录，本轮主结论不把其旧区间统计移植到四路径。

\setcounter{table}{0}
\renewcommand{\thetable}{C\arabic{table}}
\chapter{核验与复现记录}
基础与延展147个共同月收益一致；四路径初始20个季度256条名单与基础一致；P1与P2的历史候选池、三维得分一致；P2的217,083条成长率名次和263,221条PE名次与基础延展全榜一致；全部正式持仓PE为正。Python独立重算全部分段指标，并逐月重建P3、P4仓位及实际收益，与R结果在浮点精度内一致。

研究报告不附代码或原始数据。完整代码、输入版本校验、运行结果、十六组实证图和R环境信息见配套RMarkdown及其HTML；六方案结果和四条正式路径程序独立保存。原有策略程序和原始数据保持不变，修订前两份报告另作备份。

本机R保存结果并输出HTML后仍以0xC0000005退出，原因未定位。断言通过不等同于进程正常；日志及状态随附，Python独立复算和构建正常完成。

\setcounter{table}{0}
\renewcommand{\thetable}{D\arabic{table}}
\chapter{回撤恢复与完整年度的补充证据}
\section{最深月末回撤的发生与恢复}
下表从已有月收益复利计算，不生成新策略。峰值含收益开始前的净值1；同值峰顶取谷值之前最后一个月，恢复指谷值后首次达到该峰值。所列月数为两个观察月的间隔，不是日历交易日；只描述最深回撤这一段，不是所有回撤的平均时长，也不表示峰值至恢复期间始终持有同一名单。

\begin{table}[htbp]
\centering
\caption{最大月末回撤对应的时间与恢复}
\renewcommand{\arraystretch}{1.18}
\footnotesize
\setlength{\tabcolsep}{3pt}
\begin{tabularx}{\textwidth}{@{}p{1.15cm}*{6}{C}@{}}
\toprule
方案 & 峰值月份 & 谷值月份 & 回撤/\% & 至谷/月 & 首次恢复 & 至恢复/月 \\
\midrule
B & 2007-09 & 2008-10 & -70.11 & 13 & 2011-02 & 41 \\
E & 2007-09 & 2008-10 & -70.11 & 13 & 2011-02 & 41 \\
P1 & 2007-09 & 2008-10 & -68.87 & 13 & 2009-12 & 27 \\
P2 & 2007-09 & 2008-10 & -64.15 & 13 & 2009-11 & 26 \\
P3 & 2007-09 & 2008-10 & -32.74 & 13 & 2009-10 & 25 \\
P4 & 2015-05 & 2018-10 & -34.86 & 41 & 2020-06 & 61 \\
\bottomrule
\end{tabularx}
\end{table}

基准E最深回撤发生于2007年9月至2008年10月，2011年2月首次恢复此前峰值，峰值到恢复间隔41个月。P3最深回撤对应同一峰谷阶段，2009年10月恢复，间隔25个月；P4最深回撤则出现在2015年5月至2018年10月，到2020年6月恢复，共61个月。这说明较小最大回撤幅度不保证较短的恢复时间，且不同路径的最深回撤未必发生在同一事件中。

B的终点为2015年6月，其他方案为2026年9月，不能只比较“是否恢复”而忽略观察机会不同。月末统计不显示月内跌幅；原研报的日频峰谷日期不能用本表月份直接替代。每条最大回撤已与主表指标逐项复核。

\section{完整自然年度的收益分布}
\begin{table}[htbp]
\centering
\caption{2004—2025年完整年度对照}
\renewcommand{\arraystretch}{1.18}
\footnotesize
\setlength{\tabcolsep}{3pt}
\begin{tabularx}{\textwidth}{@{}p{1.15cm}*{5}{C}@{}}
\toprule
方案 & 年数 & 正收益年 & 领先E年 & 最差年 / \% & 最好年 / \% \\
\midrule
E & 22 & 16 & — & 2008 / -65.02 & 2007 / 446.28 \\
P1 & 22 & 18 & 11 & 2008 / -60.67 & 2007 / 446.28 \\
P2 & 22 & 16 & 10 & 2008 / -52.44 & 2007 / 446.28 \\
P3 & 22 & 19 & 9 & 2008 / -25.19 & 2007 / 182.09 \\
P4 & 22 & 17 & 7 & 2018 / -27.19 & 2007 / 182.09 \\
\bottomrule
\end{tabularx}
\end{table}

P1有18个正收益年度，高于E的16个，但22年中只有11年严格领先E；P3有19个正收益年度，领先E的年度仅9个。正收益年度更多与相对基准领先更多并非同一概念，绝对收益、相对收益和风险应分别阅读。P4最差完整年度为2018年，而其他路径均为2008年，和回撤时段不同的观察相呼应。

只纳入各12个月的22个年度；领先E指当年复利严格大于E，相等不计领先。E作为自身基准不报告领先次数。最好和最差仅是已观察年度极值，并非风险预测区间；跨年度次数没有控制共同市场暴露和参数选择，不当作显著性检验或未来胜率。
